% Appendix: Code Implementation and AI Use Disclosure
% MCM 2026 Problem A

\section{Appendix A: Python Implementation}
\label{app:code}

This appendix provides the core Python implementation of our smartphone battery discharge model.

\subsection{Battery Parameters Data Class}

\begin{lstlisting}[caption={Battery Parameters Definition}, label={lst:params}]
import numpy as np
from scipy.integrate import solve_ivp
from dataclasses import dataclass
from typing import Callable, Tuple, Dict

@dataclass
class BatteryParams:
    """Battery parameter data class for smartphone Li-ion cells"""
    # Capacity and voltage
    Q_nom: float = 3.0      # Nominal capacity [Ah]
    V_nom: float = 3.8      # Nominal voltage [V]
    V_min: float = 3.0      # Cutoff voltage [V]
    V_max: float = 4.35     # Maximum voltage [V]
    SOC_min: float = 0.05   # Minimum usable SOC [-]
    
    # ECM parameters (2-RC model)
    R0: float = 0.080       # Ohmic resistance [Ohm]
    R1: float = 0.020       # RC1 resistance [Ohm]
    C1: float = 2000        # RC1 capacitance [F]
    R2: float = 0.010       # RC2 resistance [Ohm]
    C2: float = 50000       # RC2 capacitance [F]
    
    # OCV-SOC polynomial coefficients
    K0: float = 3.685       # [V]
    K1: float = 0.2156      # [V]
    K2: float = 0.0351      # [V]
    K3: float = -0.0175     # [V]
    
    # Thermal parameters
    m: float = 0.050        # Battery mass [kg]
    Cp: float = 1000        # Specific heat [J/(kg*K)]
    h: float = 10           # Heat transfer coeff [W/(m^2*K)]
    A: float = 0.006        # Surface area [m^2]
    Ea_R: float = 25000     # Activation energy [J/mol]
    T_ref: float = 298.15   # Reference temperature [K]
    
    # Aging parameters
    k_cal: float = 2e-6     # Calendar aging rate [1/sqrt(s)]
    k_cyc: float = 2e-5     # Cycle aging rate [1/Ah]
    Ea_cal: float = 60000   # Calendar activation energy [J/mol]
    Ea_cyc: float = 40000   # Cycle activation energy [J/mol]
    gamma: float = 2.0      # DOD exponent [-]
    alpha: float = 0.7      # C-rate exponent [-]
    k_R: float = 2.0        # Resistance growth factor [-]

R_GAS = 8.314  # Universal gas constant [J/(mol*K)]
\end{lstlisting}

\subsection{Main Battery Model Class}

\begin{lstlisting}[caption={Smartphone Battery Model Implementation}, label={lst:model}]
class SmartphoneBatteryModel:
    """Continuous-time smartphone battery model with 2-RC ECM"""
    
    def __init__(self, params: BatteryParams = None):
        self.params = params or BatteryParams()
    
    def ocv(self, SOC: float, T: float = 298.15) -> float:
        """Open-circuit voltage model (combined polynomial-log)"""
        p = self.params
        SOC = np.clip(SOC, 0.01, 0.99)
        V_OC = p.K0 + p.K1*SOC + p.K2*np.log(SOC) + p.K3*np.log(1-SOC)
        # Temperature correction
        V_OC += -0.0003 * (T - p.T_ref)
        return V_OC
    
    def resistance(self, T: float, SOC: float, SOH: float = 1.0) -> float:
        """Temperature and SOH dependent internal resistance"""
        p = self.params
        # Arrhenius temperature correction
        R_T = p.R0 * np.exp(-p.Ea_R/R_GAS * (1/T - 1/p.T_ref))
        # SOH correction (resistance increases with aging)
        R_SOH = R_T * (1 + p.k_R * (1 - SOH))
        return R_SOH
    
    def power_to_current(self, P_mW: float, SOC: float, T: float,
                         V_RC1: float, V_RC2: float, SOH: float) -> Tuple:
        """Convert power demand to current (iterative solver)"""
        V_OC = self.ocv(SOC, T)
        R_total = self.resistance(T, SOC, SOH)
        
        # Newton iteration
        V = V_OC - V_RC1 - V_RC2
        for _ in range(5):
            I = P_mW / V / 1000  # mW to A
            V = V_OC - I * R_total - V_RC1 - V_RC2
        return I, V
    
    def battery_ode(self, t: float, x: np.ndarray,
                    power_func: Callable, T_amb: float = 298.15) -> np.ndarray:
        """
        Battery state ODE system
        State vector: x = [SOC, V_RC1, V_RC2, T]
        """
        SOC, V_RC1, V_RC2, T = x
        p = self.params
        SOH = 1.0  # Assume constant for short-term simulation
        
        # Get power demand
        P_load = power_func(t)
        
        # Solve for current and voltage
        I, V = self.power_to_current(P_load, SOC, T, V_RC1, V_RC2, SOH)
        
        # Effective capacity [As]
        Q_eff = p.Q_nom * SOH * 3600
        
        # SOC dynamics (Coulomb counting)
        dSOC = -I / Q_eff
        
        # RC dynamics with temperature-dependent time constants
        R1_T = p.R1 * np.exp(-p.Ea_R/R_GAS * (1/T - 1/p.T_ref))
        R2_T = p.R2 * np.exp(-p.Ea_R/R_GAS * (1/T - 1/p.T_ref))
        tau1, tau2 = R1_T * p.C1, R2_T * p.C2
        
        dV_RC1 = -V_RC1 / tau1 + I / p.C1
        dV_RC2 = -V_RC2 / tau2 + I / p.C2
        
        # Thermal dynamics (lumped model)
        R_total = self.resistance(T, SOC, SOH)
        Q_gen = I**2 * R_total           # Heat generation
        Q_loss = p.h * p.A * (T - T_amb) # Heat dissipation
        dT = (Q_gen - Q_loss) / (p.m * p.Cp)
        
        return np.array([dSOC, dV_RC1, dV_RC2, dT])
\end{lstlisting}

\subsection{Simulation and Prediction Methods}

\begin{lstlisting}[caption={Simulation and Time-to-Empty Prediction}, label={lst:sim}]
def simulate(self, SOC_0: float, T_0: float, power_func: Callable,
             t_span: Tuple[float, float], T_amb: float = 298.15) -> Dict:
    """
    Run battery discharge simulation
    
    Args:
        SOC_0: Initial SOC [0-1]
        T_0: Initial temperature [K]
        power_func: Power demand function P(t) [mW]
        t_span: Time interval (t_start, t_end) [s]
        T_amb: Ambient temperature [K]
    
    Returns:
        Dictionary with time, states, voltage, and current trajectories
    """
    x0 = [SOC_0, 0.0, 0.0, T_0]  # Initial state
    
    # SOC termination event
    def soc_event(t, x):
        return x[0] - self.params.SOC_min
    soc_event.terminal = True
    soc_event.direction = -1
    
    # Solve ODE
    sol = solve_ivp(
        lambda t, x: self.battery_ode(t, x, power_func, T_amb),
        t_span, x0, method='RK45', events=soc_event, max_step=10.0
    )
    
    # Compute voltage and current trajectories
    V_traj, I_traj = [], []
    for i in range(len(sol.t)):
        SOC, V_RC1, V_RC2, T = sol.y[:, i]
        P = power_func(sol.t[i])
        I, V = self.power_to_current(P, SOC, T, V_RC1, V_RC2, 1.0)
        V_traj.append(V); I_traj.append(I)
    
    return {
        't': sol.t, 'SOC': sol.y[0], 'V_RC1': sol.y[1],
        'V_RC2': sol.y[2], 'T': sol.y[3],
        'V': np.array(V_traj), 'I': np.array(I_traj),
        't_empty': sol.t[-1] if sol.t_events[0].size > 0 else None
    }

def predict_empty_time(self, SOC_0: float, P_load_mW: float,
                       SOH: float = 1.0, T: float = 298.15) -> float:
    """Quick estimation of time-to-empty [seconds]"""
    p = self.params
    Q_eff = p.Q_nom * SOH
    usable_SOC = SOC_0 - p.SOC_min
    V_avg = self.ocv((SOC_0 + p.SOC_min) / 2, T)
    
    energy_Wh = Q_eff * usable_SOC * V_avg
    P_W = P_load_mW / 1000
    
    t_hours = energy_Wh / P_W
    return t_hours * 3600  # Return in seconds
\end{lstlisting}

\subsection{Extended Kalman Filter for SOC Estimation}

\begin{lstlisting}[caption={EKF Implementation for Online SOC Estimation}, label={lst:ekf}]
class BatteryEKF:
    """Extended Kalman Filter for battery SOC estimation"""
    
    def __init__(self, params: BatteryParams, dt: float = 1.0):
        self.params = params
        self.dt = dt
        
        # State: [SOC, V_RC1, V_RC2]
        self.x = np.array([0.5, 0.0, 0.0])
        
        # Covariance matrices
        self.P = np.diag([0.01, 0.001, 0.001])  # State covariance
        self.Q = np.diag([1e-6, 1e-4, 1e-5])    # Process noise
        self.R = np.array([[0.01**2]])           # Measurement noise (V)
    
    def predict(self, I: float, T: float):
        """EKF prediction step"""
        SOC, V_RC1, V_RC2 = self.x
        p = self.params
        dt = self.dt
        
        Q_eff = p.Q_nom * 3600  # Assume SOH = 1
        tau1 = p.R1 * p.C1
        tau2 = p.R2 * p.C2
        
        # State prediction
        SOC_new = SOC - I * dt / Q_eff
        V_RC1_new = V_RC1 * np.exp(-dt/tau1) + I * p.R1 * (1 - np.exp(-dt/tau1))
        V_RC2_new = V_RC2 * np.exp(-dt/tau2) + I * p.R2 * (1 - np.exp(-dt/tau2))
        
        self.x = np.array([SOC_new, V_RC1_new, V_RC2_new])
        
        # Jacobian of state transition
        F = np.array([
            [1, 0, 0],
            [0, np.exp(-dt/tau1), 0],
            [0, 0, np.exp(-dt/tau2)]
        ])
        
        # Covariance prediction
        self.P = F @ self.P @ F.T + self.Q
    
    def update(self, V_meas: float, I: float, T: float):
        """EKF update step with voltage measurement"""
        SOC, V_RC1, V_RC2 = self.x
        p = self.params
        
        # Predicted voltage
        V_OC = p.K0 + p.K1*SOC + p.K2*np.log(max(SOC,0.01)) \
               + p.K3*np.log(max(1-SOC,0.01))
        V_pred = V_OC - I * p.R0 - V_RC1 - V_RC2
        
        # Measurement Jacobian
        dVoc_dSOC = p.K1 + p.K2/max(SOC,0.01) - p.K3/max(1-SOC,0.01)
        H = np.array([[dVoc_dSOC, -1, -1]])
        
        # Innovation
        y = V_meas - V_pred
        S = H @ self.P @ H.T + self.R
        
        # Kalman gain
        K = self.P @ H.T @ np.linalg.inv(S)
        
        # State and covariance update
        self.x = self.x + (K @ np.array([[y]])).flatten()
        self.P = (np.eye(3) - K @ H) @ self.P
        
        # Constrain SOC
        self.x[0] = np.clip(self.x[0], 0.0, 1.0)
        
        return self.x[0]  # Return estimated SOC
\end{lstlisting}

\subsection{Sobol Sensitivity Analysis}

\begin{lstlisting}[caption={Sobol Sensitivity Analysis Implementation}, label={lst:sobol}]
from SALib.sample import saltelli
from SALib.analyze import sobol

# Problem definition for sensitivity analysis
SENSITIVITY_PROBLEM = {
    'num_vars': 9,
    'names': ['brightness', 'cpu_load', 'network_power', 'gps_on',
              'T_amb', 'SOC_0', 'SOH', 'Q_nom', 'R0'],
    'bounds': [
        [0.1, 1.0],    # brightness
        [0.05, 0.95],  # cpu_load
        [50, 500],     # network_power [mW]
        [0, 1],        # gps_on
        [263, 318],    # T_amb [K]
        [0.2, 1.0],    # SOC_0
        [0.7, 1.0],    # SOH
        [2.5, 5.0],    # Q_nom [Ah]
        [0.05, 0.15]   # R0 [Ohm]
    ]
}

def run_sobol_analysis(model_func, n_samples=1024):
    """Run Sobol global sensitivity analysis"""
    # Generate samples
    param_values = saltelli.sample(SENSITIVITY_PROBLEM, n_samples)
    
    # Evaluate model
    Y = np.zeros(param_values.shape[0])
    for i, params in enumerate(param_values):
        param_dict = dict(zip(SENSITIVITY_PROBLEM['names'], params))
        Y[i] = model_func(param_dict)
    
    # Analyze
    results = sobol.analyze(SENSITIVITY_PROBLEM, Y, print_to_console=False)
    
    # Return rankings by total effect
    names = SENSITIVITY_PROBLEM['names']
    sorted_idx = np.argsort(results['ST'])[::-1]
    
    return {
        'parameters': [names[i] for i in sorted_idx],
        'S1': results['S1'][sorted_idx],
        'ST': results['ST'][sorted_idx]
    }
\end{lstlisting}

\newpage
\section{Appendix B: AI Use Disclosure}
\label{app:ai}

In accordance with MCM/ICM guidelines, we disclose the following uses of artificial intelligence tools in the preparation of this paper:

\subsection{AI Tools Used}

\begin{enumerate}
    \item \textbf{Claude (Anthropic)}: Used for literature review assistance, equation verification, and LaTeX formatting suggestions.
    
    \item \textbf{GitHub Copilot}: Used for code completion and debugging assistance in Python implementation.
\end{enumerate}

\subsection{Specific Uses}

\begin{table}[H]
\centering
\caption{Summary of AI Tool Usage}
\begin{tabular}{p{3cm}p{4cm}p{5cm}}
\toprule
\textbf{Task} & \textbf{AI Tool} & \textbf{How Used} \\
\midrule
Literature search & Claude & Suggested relevant papers on battery modeling and smartphone power consumption \\
\addlinespace
Equation checking & Claude & Verified dimensional consistency of ODE system \\
\addlinespace
Code debugging & Copilot & Auto-completed boilerplate code and suggested bug fixes \\
\addlinespace
LaTeX formatting & Claude & Assisted with TikZ diagrams and table formatting \\
\addlinespace
Writing refinement & Claude & Suggested improvements to technical writing clarity \\
\bottomrule
\end{tabular}
\end{table}

\subsection{Human Oversight}

All AI-generated content was carefully reviewed, verified, and modified by team members. The mathematical models, analysis methodology, and conclusions represent our original work. AI tools served as assistants to enhance productivity, not as primary content generators.

\subsection{Verification Process}

\begin{itemize}
    \item All equations were independently verified using dimensional analysis
    \item Code implementations were tested with known benchmark cases
    \item Literature citations were manually verified for accuracy
    \item Numerical results were cross-checked with hand calculations
\end{itemize}

\vspace{1cm}
\noindent\textit{We affirm that this paper represents our own work and that all AI assistance has been disclosed as required.}
